\chapter{Unglaube und Aberglaube}
\quellseite{8}

„Ich sehe euch, daß ihr in allen Stücken allzu abergläubig seid“, sagte St. Paulus zu den stolzen, gebildeten Athenern, Apost. 17, 22. Zu wie vielen kann das gesagt werden! Etwas muß der Mensch freilich glauben. Aber wo er nicht an Gott glaubt, da wuchert der Aberglaube. Mose und den Propheten glaubt der reiche Mann nicht, aber das Lazarusgespenst soll Wunder wirken. Voltaire, ein Meister im Spott über die Bibel, kam immer betrübt nach Haus, so oft er die Raben auf dem Felde hatte zur Linken krächzen gehört. Herzog Philipp von Orleans, der ein großer Freigeist war und in der französischen Revolution mitgeholfen hatte, den lieben Gott abzusetzen, und viele zum Tode gebracht hatte, kam zuletzt ins Gefängniß. Dort ließ er sich mit banger Seele aus dem Kaffeesatze lesen, ob er
freigesprochen
\quellseite{9}
oder hingerichtet werden würde. — Wie anders dagegen das Bild des Landgrafen Wilhelm von Hessen, der, als ihm das Buch eines Sterndeuters gezeigt wurde, worin des Landgrafen Sterbetag bezeichnet war, ruhig an den Rand jener Stelle schrieb: Psalm 31, 16.: „Meine Zeit stehet in Deinen Händen!“
